% Folder 142, batch 04 — archivists' pages 61–80.
% Pass: Opus 5.5 (claude-opus-5-5), 2026-10-02 — first pass, unchecked against the pages by a human.
%
% Pages 61–63 close a run on the sphere with marked points (formulas (41)–(45),
% author's p. 12 on page 62). Page 64 is blank. Page 65 is a cover in his hand;
% its words are the section heading. Pages 66–80 open a new run, paginated by
% him 1–8 on the even archivists' pages; it continues past page 80.
\documentclass[11pt,a4paper]{article}
\input{../preamble/grothendieck}

\folder{142}
\batch{4}
\pages{61}{80}
\foldertitle{Action de Galois sur Teichmüller : notes manuscrites (s.d.).}
\dating{[à partir de 1978]}
\watermark{Édition de démonstration}

\begin{document}

\page{61}
\note{la page commence au milieu d'une phrase : l'argument vient de la page précédente, hors de ce lot ; les numéros de formules (41)–(45) et les notations $\widetilde{\Sigma}$, $\Pi_{0,3}$, $\Psi$, $\varphi_r$ y renvoient}

\drawing{En haut à gauche, une sphère : pôle nord $P_+$, pôle sud $P_-$ ; sur l'équateur, de gauche à droite, $Q_1$, $R_0^-$, $R_0^+$, $Q_\infty$, $R_1^-$ ; deux petits cercles sur l'équateur, l'un entre $R_0^-$ et $R_0^+$ (avec $S_0^+$ au-dessus, $S_0^-$ au-dessous), l'autre près de $R_1^-$ (avec $S_1^+$) ; des méridiens joignent $P_+$ à $P_-$ ; un fuseau hachuré sous $P_+$, marqué d'une flèche circulaire d'orientation.}

correspondant, sur la sphère standard, à la région \uncertain{circulaire}, délimitée par le contour
\[
P_+\,Q_1\,R_0^-\,S_0^+\,R_0^+\,S_0^-\,P_-\,Q_\infty\,R_1^-\,S_1^+\,P_+ .
\]
\note{le $S_0^+$ du contour est écrit sur une autre lettre}
Les régions $\widetilde{\Sigma}_r$ sont permutées de façon simplement transitive par $\Gamma = \mathfrak{S}_3 \times (1,\tau)$. Elles sont simplement connexes, et permettent donc de définir des
\[
\Pi_r = \Pi_i^{\omega,\varepsilon} = \pi_1(\widetilde{\Sigma};\widetilde{\Sigma}_r) = \pi_1(\Sigma^{*};\Sigma^{*}_r)
\]
\note{numéro de formule surchargé devant cette ligne, illisible ; l'exposant $\omega,\varepsilon$ pourrait aussi se lire $\omega,i$, comme en (42)}
\struck{et des iso} \struck{groupes} mais qui sont reliés par un syst. transitif d'iso. à l'aide des \uncertain{opérations} de $\Gamma$, de sorte qu'on peut identifier tous ces groupes entre eux. En fait, on trouve des isom. d'épinglage
\[
\text{(42)}\qquad \Psi_r : \Pi_{0,3} \longrightarrow \Pi_r = \Pi_i^{\omega,i} = \pi_1(\widetilde{\Sigma};\widetilde{\Sigma}_r)
\]
($\Psi_r$ \ill{}) \struck{en identi} en composant

\begin{tikzcd}
\Pi_{0,3} \arrow[r, "\Psi_r = \Psi_\varepsilon^{\omega,i}"] & \pi_1(\widetilde{\Sigma},P_\varepsilon) \simeq \pi_1(\Sigma^{*},P_\varepsilon) \arrow[r, leftrightarrow, "\varphi_r^{-1}"] & \Pi_r
\end{tikzcd}

\note{la flèche $\varphi_r^{-1}$ porte une pointe à chaque bout ; l'isomorphisme $\simeq \pi_1(\Sigma^{*},P_\varepsilon)$ est écrit sous le terme du milieu}

où \struck{$\varphi$} $\varphi_r : \Pi_r \to \Pi_\varepsilon = \pi_1(\widetilde{\Sigma},P_\varepsilon)$ est \struck{la} l'isomorphisme de transition « correspondant » à la région 1-connexe $\Sigma_r$, et au point $P_\varepsilon$ de $\Sigma_r$.

\page{62}
\note{p. 12 de l'auteur}
La donnée d'un $R_i^{\omega}$\note{les indices de ce premier $R$ sont surchargés ; lu d'après la suite}, \add{qu'on va écrire $R^{\omega}_{\omega i}$, équivalant à} celle d'une région \add{($r \neq (i,\omega)$)} de $J$, et\note{la relation $r \neq (i,\omega)$, écrite au-dessus de la ligne, est douteuse} définit \emph{deux} \struck{triang} régions $\widetilde{D}_r$ qui le contiennent, correspondant aux deux régions de ($J$, \ill{}) qui prolongent $r$ — et il y a donc deux régions correspondantes $\widetilde{\Sigma}_r$, qui sont les seules qui contiennent $R_i^{\omega}$ comme \emph{point intérieur} (NB. dans la région $\widetilde{\Sigma}_r$ standard, les deux autres points $R$ : par $R_0^+$, qui est intérieur, sont $R_0^-$ et $R_1^-$, qui ne le sont pas). Mais si la sphère $\Sigma$ est orientée i.e. si on s'est donné un élément marqué \add{$\varpi_+$} de $\underline{\varpi} = \underline{\omega} \wedge \underline{\varepsilon}$, on aura un isom $\underline{\omega} \simeq \underline{\varepsilon}$, donc le point $R_i^{\omega}$ détermine un $\varepsilon$ \add{$\in \underline{\varepsilon}$} privilégié, par la condition que $\varepsilon$ corresponde à $\omega$ par $\underline{\omega} \simeq \underline{\varepsilon}$, ou encore que $\omega \wedge \varepsilon = \varpi_+$, ce qui signifie aussi que l'orientation $P_+S_0^+R_0^+Q_\infty P_+$ de la région \struck{$\widetilde{\Sigma}$} $\widetilde{D}_r$ est \struck{\ill{}} $\varpi_+$. \uncertain{Donc} on peut alors identifier canoniquement, via $\Psi_r$, $\pi_1(\widetilde{\Sigma},R_i^{\omega})$ à $\Pi_{0,3}$, par un iso (qui dépend du choix de l'orientation $\varpi_+$ de $\Sigma$)
\[
\text{(43)}\qquad \Psi_i^{\omega} : \Pi_{0,3} \xrightarrow{\ \sim\ } \pi_1(\widetilde{\Sigma},R_i^{\omega})
\]

\page{63}
de sorte qu'on aura \struck{compatibilité} \add{\uncertain{une compatibilité}} des

\[\text{(44)}\]

\begin{tikzcd}[column sep=small]
& \Pi_{0,3} \arrow[dl, "\Psi_i^{\omega}"'] \arrow[d, "\Psi_i^{\omega}"] \arrow[dr, "\Psi_\omega^{i}"] & \\
\pi_1(\widetilde{\Sigma},R^{\omega}_{\omega i}) \arrow[r, leftrightarrow, "b_i^{\omega'}(\cdot)"'] & \pi_1(\widetilde{\Sigma},Q_i) \arrow[r, leftrightarrow, "\varphi_{i\omega} = c_i^{\omega}(\cdot)"'] & \pi_1(\widetilde{\Sigma},P_\omega) \\
& = \Pi_i & = \Pi_\omega
\end{tikzcd}

\note{la flèche de gauche est tracée deux fois, sur un $\Pi_{0,3}$ biffé ; l'indice de l'étiquette de la flèche verticale est corrigé en $i$}

où on a écrit $P_\omega$ au lieu de $P_\varepsilon$ (en identifiant $\underline{\omega}$ et $\underline{\varepsilon}$ grâce à $\varpi_+$) et $\Psi_i^{\omega}$ au lieu de $\Psi_i^{\omega,\omega}$, $\Psi_\omega^{i}$ au lieu de $\Psi_\omega^{i,\omega}$.

Le passage à expliciter est celui de \struck{$R_{\omega i}$}
\[
\Pi_i^{\omega} = \pi_1(\widetilde{\Sigma},R^{\omega}_{\omega i}) \quad\text{à}\quad \Pi_i^{\omega'} = \pi_1(\widetilde{\Sigma},R^{\omega'}_{\omega' i}),
\]
chacun identifié à $\Pi_{0,3}$ par $\Psi_i^{\omega}$, resp. $\Psi_i^{\omega'}$\note{écrit sous chaque terme : une flèche verticale $\uparrow\wr$ marquée $\Psi_i^{\omega}$, resp. $\Psi_i^{\omega'}$, issue de $\Pi_{0,3}$}, \add{via $\widetilde{b}_{i\omega}$.} On trouve par la commutativité (44) et par la deuxième formule (39) qu'on a

\drawing{En marge gauche : deux petits cercles, aux points $R^{\omega}_{\omega i}$ et $R^{\omega'}_{\omega' i}$, joints par un chemin fléché $\widetilde{b}_i^{\omega}$ qui passe par $Q_i$.}

\[
\text{(45)}\qquad \widetilde{b}_i^{\omega}\bigl(\Psi_i^{\omega}(x)\bigr) = \Psi_i^{\omega'}\bigl(\sigma_\infty(x)\bigr) .
\]
\note{devant (45), un $\widetilde{b}_{\omega i}$ biffé}

\section{Généralités sur les actions de groupes de Galois sur les compactifiés profinis de groupoïdes fondamentaux de variétés alg. complexes}
\note{titre de sa main, sur le feuillet de couverture (p. 65), avec au crayon « (38 pages) »}

\page{66}
\note{p. 1 de l'auteur}
Soit $K$ un sous-corps du corps des complexes $\mathbb{C}$ — des cas intéressants seront celui où $K = \mathbb{Q}$, ou $K$ extension finie de $\mathbb{Q}$, ou $K = \widetilde{\mathbb{Q}}$ extension algébrique maximale de $\mathbb{Q}$ dans $\mathbb{C}$ — je serai intéressé en tous cas avant tout \uncertain{dans} cas où $K$ est de t.f. sur $\mathbb{Q}$. (Le cas de $\widetilde{\mathbb{Q}}$ est un cas « limite » de tels cas, pour ne pas être gêné par les racines de 1…).

On s'intéresse aux schémas $X$ de t.f. sur $K$, et aux espaces \add{analytiques} $X(\mathbb{C})$ correspondants. Comme on s'intéresse aux invariants « topologiques » de $X$, $X(\mathbb{C})$, on traitera les morphismes $X' \to X$ qui sont des homéomorphismes universels comme des iso — en particulier les éléments nilpotents ne comptent pas — une immersion fermée surjective sera traitée comme un iso.

Un langage commode est de considérer $X$ comme donné par l'espace analytique (ou même seulement topologique) $X(\mathbb{C})$, plus une structure \emph{supplémentaire} dessus : la structure de « var. alg. sur $K$ » ; \uncertain{à savoir} la top. de Zariski, le faisceau des fonctions $K$-polynomiales etc. On écrira par la suite $X$ au lieu de $X(\mathbb{C})$.

On distingue sur $X$ le groupoïde fondamental $\Pi_1(X)$ au sens topologique, et son compactifié profini $\widehat{\Pi}_1(X)$, qui \struck{est aussi} \add{peut} aussi être défini directement en termes du $\mathbb{C}$-schéma $X$. On a un foncteur « d'inclusion »

\page{67}
\[
\text{(1)}\qquad \Pi_1(X) \longrightarrow \widehat{\Pi}_1(X)
\]
qui induit une bijection sur les ensembles des objets \add{(ou « sommets »)} — s'identifiant à $X = X(\mathbb{C})$ — et qui \struck{défi} pour tout $x \in X = X(\mathbb{C})$, donne l'homom de groupes
\[
\text{(2)}\qquad
\underbrace{\operatorname{Aut}_{\Pi_1(X)}(x)}_{\overset{\text{déf}}{=}\ \pi_1(X,x)} \longrightarrow \underbrace{\operatorname{Aut}_{\widehat{\Pi}_1(X)}(x)}_{\overset{\text{déf}}{=}\ \widehat{\pi}_1(X_{\mathrm{sch}},x)}
\]
permettant d'identifier le $\pi_1$ schématique au compactifié profini du $\pi_1$ « topologique » ou « transcendant » :
\[
\text{(3)}\qquad \pi_1(X_{\mathrm{sch}},x) \simeq \widehat{\pi}_1(X,x) .
\]
De même, \add{grâce à (3)}, pour $x,y \in X$, l'application
\[
\operatorname{Hom}_{\Pi_1(X)}(x,y) \longrightarrow \operatorname{Hom}_{\widehat{\Pi}_1(X)}(x,y)
\]
permet d'identifier le deuxième ens, en tant que bitorseur sous $\bigl(\widehat{\pi}_1(X,y),\widehat{\pi}_1(X,x)\bigr)$, au bitorseur déduit \struck{Hom} du premier ens par extension des groupes d'opérateurs par
\[
\pi_1(X,y) \to \widehat{\pi}_1(X,y), \qquad \pi_1(X,x) \to \widehat{\pi}_1(X,x)
\]
(\ill{} compte d'homomorphismes satisfaisant à la condition \add{bien} connue requise pour que cette extension ait un sens…).

\marginal{\ill{} — note écrite en oblique dans la marge gauche, près de la formule (3) ; on y lit « On identifiera… », $\Pi_1$, $(x,y)$, (9)}
\marginal{\ill{} — seconde note oblique en bas de la marge gauche, avec une flèche vers le texte ; on y lit $(G,H)$, $(G',H')$, $G \to G'$, $H \to H'$}

Soit
\[
\text{(5)}\qquad \Gamma = \Gamma_K \overset{\text{déf}}{=} \operatorname{Aut}_{K\text{-cat}}(\mathbb{C}),
\]
\note{l'indice d'un premier $\Gamma$ est biffé ; aucune formule (4) n'est numérotée sur la page}

\page{68}
\note{p. 2 de l'auteur}
dont $\Gamma$ opère non seulement sur \struck{$X$} l'ens. $X$, mais aussi sur le groupoïde $\widehat{\Pi}_1(X)$, à cause du fait que celui-ci admet une définition purement schématique, ne faisant pas intervenir la topologie de $\mathbb{C}$. Bien entendu, si $\lambda : x \to y$ dans $\Pi_1$, alors pour $u \in \Gamma$, $u(\lambda)$ n'est pas en général une flèche de $\Pi_1$ (de $u(x)$ vers $u(y)$) mais seulement de $\widehat{\Pi}_1$.

\struck{Notation} Notations ${}^{u}x$ ($x \in X$, $u \in \Gamma$) et
\[
{}^{u}\lambda : {}^{u}x \longrightarrow {}^{u}y \quad\text{si}\quad \lambda : x \to y \ \ (\text{dans } \widehat{\Pi}_1(X)) .
\]
On a
\[
\text{(6)}\qquad
\begin{cases}
u(\mathrm{id}_x) = \mathrm{id}_{{}^{u}x} & (x \in X)\\
u(\mu\lambda) = {}^{u}\mu\,{}^{u}\lambda & \Bigl(\begin{array}{l} x \xrightarrow{\lambda} y \xrightarrow{\mu} z \\ {}^{u}x \xrightarrow{{}^{u}\lambda} {}^{u}y \xrightarrow{{}^{u}\mu} {}^{u}z \end{array}\Bigr)
\end{cases}
\]

Notre objet est d'essayer de « déterminer » autant que faire se peut cette action de $\Gamma$. À vrai dire, si $\overline{K}$ est la clôture algébrique de $K$ dans $\mathbb{C}$, on s'intéressera surtout aux sous-groupoïdes de $\Pi_1(X)$ et $\widehat{\Pi}_1(X)$ induits sur l'ens d'objets $X_{\overline{K}}$, qu'on désignera par $\Pi_1(X_{\overline{K}})$, $\widehat{\Pi}_1(X_{\overline{K}})$ — le deuxième est stable par l'action de $\Gamma$, et si on pose
\[
\text{(7)}\qquad \overline{\Gamma} = \operatorname{Aut}_{K\text{-cat}}(\overline{K}) = \operatorname{Gal}_{\overline{K}/K} ,
\]
de sorte que $\overline{\Gamma}$ est un quotient de $\Gamma$, l'action de $\Gamma$ sur $\Pi_1(X_{\overline{K}})$ se fait à travers $\overline{\Gamma}$.
\note{la barre sur le $\Gamma$ de (7) et sur celui de la dernière ligne est peu marquée ; l'action porte sur $\widehat{\Pi}_1(X_{\overline{K}})$ d'après ce qui précède, mais la page écrit ici $\Pi_1(X_{\overline{K}})$ sans chapeau}

\page{69}
Pour décrire les actions d'un $u \in \overline{\Gamma}$ (resp. $u \in \Gamma$) sur $\widehat{\Pi}_1(X)$ (resp. \struck{$\Pi$} $\widehat{\Pi}_1(X_{\overline{K}})$), on prendra comme « générateurs » du groupoïde $\widehat{\Pi}_1$ des chemins provenant de $\Pi_1$. On aura, pour $l : x \to y$ dans $\Pi_1$, ${}^{u}l : {}^{u}x \to {}^{u}y$. Si \struck{\ill{}} par la suite, pour simplifier, on se borne à des points $x, y, \ldots$ rationnels sur $K$ — i.e. qu'on travaille dans le sous-groupoïde $\widehat{\Pi}_1(X_K)$ formé des objets qui sont invariants sous $\Gamma$ — alors on aura
\[
\text{(*)}\qquad {}^{u}l : {}^{u}x = x \longrightarrow {}^{u}y = y ,
\]
et on peut écrire, \struck{\ill{}}…
\[
\text{(8)}\qquad
\begin{cases}
{}^{u}l = g_l(u)\circ l = l \circ f_l(u) \\
g_l(u) \in \widehat{\pi}_1(X,y) \\
f_l(u) \in \widehat{\pi}_1(X,x)
\end{cases}
\]
\note{à droite de la première ligne, un « $g_l(u) = {}^{u}l \circ l^{-1}$ » est biffé}
où $f_l(u)$, $g_l(u)$, ${}^{u}l$ se déterminent mutuellement, \struck{$f_l$ et $g_l$} par les formules (8) et
\[
\text{(9)}\qquad
\begin{cases}
f_l(u) = l^{-1}\cdot{}^{u}l , & g_l(u) = {}^{u}l\cdot l^{-1} \\
g_l(u) = l\bigl(f_l(u)\bigr) , & f_l(u) = l^{-1}\bigl(g_l(u)\bigr)
\end{cases}
\]

\marginal{On écrira \uncertain{aussi} $g_u(l)$, $f_u(l)$ au lieu de $g_l(u)$, $f_l(u)$ — note oblique dans la marge gauche, en face de (8), lecture incertaine}

On décrira l'action d'un $u \in \Gamma$ \add{sur un sous-groupoïde plein de $\Pi_1(X)$} à l'aide des éléments $f_l(u)$, ou $g_l(u)$, \add{où $l$} \uncertain{parcourt} un syst. de générateurs convenable du groupoïde plein correspondant de $\Pi_1(X)$. Dans le cas \uncertain{d'un seul}
\note{la phrase se poursuit sur la page suivante}

\page{70}
\note{p. 3 de l'auteur}
sous-groupoïde induit sur une partie finie \add{$I$} de $X_K$, on pourra trouver un ens fini de tels générateurs — il suffit p.ex. de choisir un \struck{\ill{}} $x_0 \in I$, de choisir pour tout $x \in I \setminus x_0$ un chemin $l_x$ de $x_0$ à $x$, et de plus un syst. \add{fini} de générateurs $l_i$ de $\pi_1(X,x_0)$.

\marginal{trait vertical double dans la marge, le long du paragraphe qui suit}
Je \uncertain{présume} que sous des conditions assez générales (p.ex. si $X$ est une courbe alg. \uncertain{anabélienne}) la connaissance de l'opération $u$ sur $\widehat{\Pi}_1(X_K, I)$ \add{\uncertain{pour un}} ($u \in \Gamma$) \struck{suffit à} détermine $u$ — et la question se pose alors de voir comment la connaissance des invariants $f_{l_\alpha}(u)$ (pour un syst. générateur de flèches $l_\alpha$ de $\Pi_1(X, I)$) \add{\uncertain{dis}} permet de retrouver les invariants $f_l(u)$ \uncertain{associés} à d'autres $K$-schémas $Y$…

Quand il s'agit du même schéma $X$, et qu'on s'intéresse à d'autres chemins $l : x \to y$ ($x, y \in I$) que les générateurs choisis là, la formule (6) donne en principe un moyen de \uncertain{reconstituer} tous les $g_l(u)$ pour tout $l$, puisqu'il permet de retrouver les ${}^{u}l$, et que ${}^{u}l$ est multiplicatif en $l$. En termes des $g_l(u)$, \struck{$f_l(u)$,} la formule de multiplicativité (6) s'écrit, \struck{pour} une situation
\[
x \xrightarrow{\ l'\ } y \xrightarrow{\ l\ } z
\]
\[
\text{(10)}\qquad g_{ll'}(u) = g_l(u)\cdot l\bigl(g_{l'}(u)\bigr) \quad\text{soit}\quad g_u(ll') = g_u(l)\cdot l\bigl(g_u(l')\bigr)
\]
\note{le premier membre de (10) est surchargé ; la flèche $y \to z$ est marquée d'un $l$ écrit sur un autre signe}

\marginal{NB pour la \ill{}, on ne s'intéressera qu'aux $x \in X_K$ qui sont dans $I$, \ill{} syst. \ill{} finis de \ill{} (6) — note oblique dans la marge gauche, lecture très incertaine}

\page{71}
et en termes des $f$ \ill{} la formule, essent. identique mais d'allure un peu moins sympathique
\[
\text{(11)}\qquad f_{ll'}(u) = l'^{-1}\bigl(f_l(u)\bigr)\cdot f_{l'}(u) \quad\text{i.e.}\quad f_u(ll') = l'^{-1}f_u(l)\cdot f_u(l')
\]
\note{le second membre de droite porte $l^{-1}$ ou $l'^{-1}$ : l'apostrophe n'est pas sûre}

\note{le paragraphe qui suit, jusqu'à « par composition », est barré de trois longs traits obliques}
\struck{La formule (10) peut s'exprimer en disant que pour $l$ fixé, $u \in \Gamma$ variable,}
\[
\text{(12)}\qquad \text{\struck{$g_l : \Gamma \longrightarrow \widehat{\pi}_1(X,y)$}}
\]
\struck{est un 1-cocycle sur $\Gamma$, à valeurs dans $\widehat{\pi}_1(X,y)$, relativement à l'action naturelle de $\Gamma$ sur $\widehat{\pi}_1(X,y)$. Bien entendu, les $\widehat{\pi}_1(X,y)$-torseurs à gauche,} \add{\struck{sous $\Gamma$,}} \struck{correspondant à ces 1-cocycles, n'est autre que $\operatorname{Hom}_{\widehat{\Pi}_1(X)}(x,y)$, sur lequel $\widehat{\pi}_1(X,y) = \operatorname{Hom}_{\widehat{\Pi}_1(X)}(y,y)$ opère par composition.}

Pour un composé d'un \struck{fini} nombre fini composable de chemins
\[
x_0 \xrightarrow{\ l_1\ } x_1 \xrightarrow{\ l_2\ } \cdots \; x_{n-1} \xrightarrow{\ l_n\ } x_n
\]
la formule (10) prend la forme
\[
\text{(12)}\qquad g_{l_n l_{n-1}\cdots l_1}(u) = g_{l_n}(u)\cdot l_n\bigl(g_{l_{n-1}}(u)\bigr)\cdot l_n l_{n-1}\bigl(g_{l_{n-2}}(u)\bigr)\cdots
\]
\[
\cdots (l_n\cdots l_2)\,g_{l_1}(u) .
\]
\note{le numéro (12) est surchargé — il a pu être corrigé depuis (13) une fois le (12) du passage barré abandonné ; le dernier facteur est ajouté au-dessus de la ligne}

Formules les plus \struck{évidentes} \add{simples} possibles assurément de celles qu'on pourrait imaginer qui ont un sens. On ne \struck{développera} \add{\uncertain{voit}} comment \uncertain{retrouver} que \struck{même} d'écrire cette formule, ainsi que \struck{(11)}, la formule analogue en $f$ et les formules
\note{la phrase se poursuit sur la page suivante}
\page{72}
\note{p. 4 de l'auteur}
(9) de passage de $f$ à $g$, sous une forme considérablement plus simple.

Soit $T$ un espace topologique \struck{simplement} 1-connexe, et
\[
\text{(13)}\qquad p : T \longrightarrow X
\]
une application continue \add{(la conjugaison sera la plus intéressante)}. Soit $T_K = p^{-1}(X_K)$. \struck{Considérons $x, y \in T_K$} \struck{$\Pi_1(T)$, \ill{}}
\note{les numéros de formules (13) à (17) de cette page sont tous surchargés — il semble qu'il ait renuméroté d'une unité}
\[
\text{(14)}\qquad \text{\struck{$\Pi_1(X;$}}\ \pi_1(X;T) = \pi
\]
alors pour tout $x \in T$ on a un isom canonique
\[
\text{(15)}\qquad \varphi_x : \pi \xrightarrow{\ \sim\ } \pi_1\bigl(X,p(x)\bigr)
\]
\add{de sorte que les $\pi_1(X,p(x))$ sont munis d'un syst. transitif d'isom} et pour $x, y \in T$, l'unique \struck{\ill{}} \add{flèche} de $\operatorname{Hom}$ \add{\uncertain{ou} \ill{}} de $x$ à $y$ dans $\Pi_1(T)$ donne une flèche canonique
\[
\text{(16)}\qquad l_{y,x} : p(x) \longrightarrow p(y) \quad\text{dans } \Pi_1(X)
\]
(NB cette flèche dépend du choix de $x, y$ au-dessus de $px$, $py$ — mais seulement si $p$ non injective).

\marginal{On dira que $T$ \ill{} pour \ill{} (15), \ill{} — note oblique dans la marge gauche, en face de (15)}

Moyennant les iso (15), on peut considérer \add{\uncertain{dès lors}} les $f_{l_{y,x}}(u)$ et $g_{l_{y,x}}(u)$ (si $x, y \in T_K$ i.e. $p(x), p(y) \in X_K$) comme des éléments du même groupe $\widehat{\pi}$ \add{(\uncertain{$\widehat{\phantom{x}}$})}, \struck{\ill{}} et en fait ces deux éléments \uncertain{sont égaux} au \ill{} et au même élément, soit
\[
\text{(17)}\qquad g^{T/X}_{y,x}(u) = g^{p}_{y,x}(u) = \varphi_y^{-1}\bigl(g_{l_{y,x}}(u)\bigr) = \varphi_x^{-1}\bigl(f_{l_{y,x}}(u)\bigr) \qquad \bigl(\in \widehat{\pi},\ \pi = \pi_1(X;T)\bigr) .
\]
\note{sous le premier membre de (17), un $\cap$ suivi de « $\pi = \pi_1(X;T)$ » : lu comme l'appartenance à $\widehat{\pi}$}

\page{73}
Les formules (10) (11) se réduisent alors à la formule
\[
\text{(18)}\qquad g_{zy}(u)\,g_{yx}(u) = g_{zx}(u) \qquad x,y,z \in T_K
\]
\note{le numéro (18) est surchargé ; dans (18)–(20), l'exposant $T$ des $g$ est sous-entendu}
\marginal{On a ainsi \uncertain{\ill{}} l'exposant $T$ : $g^{T/X}_{yx}$, qui \uncertain{est} \ill{} si $T \neq$ \ill{}… — note oblique dans la marge gauche, en face de (18)}
la formule (12) à la formule itérée de celle-ci
\[
\text{(19)}\qquad g_{x_n,x_0} = g_{x_n,x_{n-1}}\,g_{x_{n-1},x_{n-2}}\cdots g_{x_1,x_0}
\]
— quant aux formules de passage (9), elles \struck{\ill{}} disparaissent !

Bien entendu, les invariants $g_l(u)$ ou $f_l(u)$ \struck{peuvent} être considérés comme déduits par des invariants précédents du type $g^{T}_{y,x}$, en prenant p.ex. $T = [0,1]$, $p$ une application $[0,1] \to X$ \struck{définissant} \add{\uncertain{définissant}} \struck{un} \uncertain{un} chemin $l$ (plutôt, à homotopie près de chemins…) de sorte qu'on aura
\[
\text{(20)}\qquad f_l(u) = \varphi_x\bigl(g^{T}_{y_0,x_0}(u)\bigr), \qquad g_l(u) = \varphi_y\bigl(g^{T}_{y_0,x_0}(u)\bigr)
\]
\[
\text{\struck{$x_0 = 0 \in T$, $y_0$}} \quad x_0 = 0,\ y_0 = 1 \in T = [0,1] .
\]
\marginal{Autre exemple « illuminant » : $x_0 \in X_K$ fixé, et $T$ « l'espace des chemins » d'origine $x_0$ ($\ill{}$ les $y \in X_K$, $\ill{}$ « $x_0 \to y$ », \ill{} $\gamma_{y_0} \in \Pi_1(X;x_0)$) — note oblique dans la marge gauche, en face de (20), lecture incertaine}

\struck{La relation des 1-cocycles} \struck{Bien entendu} \struck{De fait} Pour comparer les invariants $g^{T}_{x,y}$ associés à différents espaces 1-connexes au-dessus de $X$, il suffit de regarder le cas où on a un diag. commut.

\[\text{(21)}\]

\begin{tikzcd}[column sep=small]
T \arrow[rr, "f"] \arrow[dr, "p"'] & & T' \arrow[dl, "p'"] \\
& X &
\end{tikzcd}

d'où
\[
f_K : T_K \longrightarrow T'_K ,
\]
\note{la phrase se poursuit sur la page suivante}

\page{74}
\note{p. 5 de l'auteur}
et un homomorphisme
\[
\text{(22)}\qquad \pi = \pi_1(X;T/X) \xrightarrow{\ \psi = \pi_1(f)\ } \pi' = \pi_1(X,T'/X)
\]
Ceci dit, si $x, y \in T_K$, d'où $x' = f(x)$, $y' = f(y) \in T'_K$, on aura
\[
\text{(23)}\qquad g^{T'}_{y',x'}(u) = \psi\bigl(g^{T}_{y,x}(u)\bigr)
\]
\[
\bigl(x' = f(x),\ y' = f(y),\ \psi = \pi_1(f)\bigr)
\]
Conséquence pratique : supposons $X = \bigcup_{\alpha \in A} T_\alpha$, les $T_\alpha$ 1-connexes, d'où des \struck{fonctions} \add{groupes}
\[
\text{\struck{$g^{\alpha} : T_{\alpha K} \times T_{\alpha K} \to$}}
\]
\[
\text{(24)}\qquad \pi_\alpha \overset{\text{déf}}{=} \pi_1(X, T_\alpha)
\]
et des fonctions
\[
\text{(25)}\qquad
\left\{
\begin{array}{l}
T_{\alpha K} \times T_{\alpha K} \times \Gamma \longrightarrow \widehat{\pi}_\alpha \\
(x,y,u) \longmapsto g^{T_\alpha}_{y,x}(u) \overset{\text{déf}}{=} g^{\alpha}_{y,x}(u)
\end{array}
\right.
\]
\note{le but de (25) est écrit $\pi_\alpha$ sans chapeau}
satisfaisant la formule de multiplicativité (19). Pour relier les fonctions entre elles, \struck{il faut} on \struck{considère} \add{note que} \struck{\ill{}} pour $T_\alpha \subset T_\beta$, on a
\[
\text{(26)}\qquad \text{\struck{$\ill{}$}}\ \ \varphi_{\beta\alpha} : \pi_\beta \longrightarrow \pi_\alpha
\]
et les relations
\[
\text{(27)}\qquad g^{\beta}_{y,x}(u) = \varphi_{\beta\alpha}\bigl(g^{\alpha}_{y,x}(u)\bigr) \qquad x, y \in T_{\alpha K},\ u \in \Gamma
\]
\note{$\varphi_{\beta\alpha}$ est écrit $\pi_\beta \to \pi_\alpha$ en (26) mais appliqué à $g^{\alpha}$ pour donner $g^{\beta}$ en (27) : le sens de la flèche ne s'accorde pas, la page le laisse ainsi}

Supposons que l'ens d'indices $I$, muni de la relation d'ordre induite par l'inclusion $T_\alpha \subset T_\beta$, soit tel que \struck{$\pi_0(I) \to \pi_0(X)$ soit bijectif}
\[
\text{(28)}\qquad \forall \alpha, \beta \in I,\ \ T_\alpha \cap T_\beta = \bigcup_{\gamma \leq \alpha,\beta} T_\gamma
\]
\note{sous (28), un « $\pi_0(I) \xrightarrow{\sim} \pi_0(X)$ » est biffé ; la phrase se poursuit sur la page suivante}

\page{75}
\note{le haut de la page, jusqu'à « Supposons de plus », est barré de trois traits obliques}
\struck{(cela assure la surjectivité \ill{}) que, si $T_\alpha$, $T_\beta$ sont dans une même comp. connexe de $X$, il existe une « chaîne » d'\ill{}}
\[
\text{\struck{$i_0 = i,\ i_1, \ldots, i_n = j$}}
\]
\struck{telle que pour deux indices consécutifs $k = i_\nu$, $k' = i_{\nu+1}$, on ait $T_k \subset T_{k'}$ ou $T_{k'} \subset T_k$.}
Supposons de plus que les intérieurs des $T_i$ recouvrent $X$, \struck{(\ill{})} \add{que les $T_i$ soient \ill{}} \ill{} je dis que la connaissance des
\[
g^{\alpha}_{y,x}(u) \in \widehat{\pi}_\alpha \qquad (\alpha \in A,\ y, x \in T_\alpha)
\]
\add{\uncertain{\ill{}}} (pour $u$ fixé) permet de retrouver tous les invariants $g_l(u)$ (ou $f_l(u)$). En effet, il suffit de voir que les flèches $l^{\alpha}_{y,x} = p_\alpha(l_{y,x})$ ($\alpha \in I$, $x, y \in T_\alpha$, $p_\alpha : T_\alpha \hookrightarrow X$ l'inclusion) \struck{\ill{}} forment un syst. de générateurs du groupoïde $\Pi_1(X_K)$ — \add{tous} les objets seulement sont obtenus \struck{\ill{}} en tous cas, et on trouve une application bijective sur les $\pi_0$ (du groupoïde engendré par les gén. $l^{\alpha}_{y,x}$ \add{($\alpha \in I$, $x, y \in T_{\alpha K}$)} et les relations $l^{\beta}_{y,x} = \varphi_{\beta\alpha}(l^{\alpha}_{y,x})$, et de $\Pi_1(X_K)$), il faut voir qu'on a une application surjective sur les $\pi_1$. Quand les intérieurs des \struck{$T_\alpha$} \add{$T_i$} recouvrent $X$, il suffit de noter que la connaissance des $g^{\alpha}_{y,x}(u)$ nous donne les $g_{y,x}(u)$ (sans exposant $\alpha$ !) pour $x, y$ points « voisins » sur $X$,
\note{la phrase se poursuit sur la page suivante}

\marginal{Supposons de plus que les intérieurs des $T_\alpha$ recouvrent $X$, et que les $T_\alpha$ soient \ill{} (\ill{} $A$ fini). Alors $\pi_0(A) \xrightarrow{\sim} \pi_0(X)$. Mieux — note oblique dans la marge gauche, en haut de page, lecture incertaine}

\page{76}
\note{p. 6 de l'auteur}
et tout chemin effectif $x \to y$ (donné par une appl continue $[0,1] \to X$) se décompose en \add{composé de} chemins \add{« directs »} $x_0 = x \to x_1 \to x_2 \to \cdots \to x_n = y$ entre \struck{\ill{}} \add{points} « proches ». Dans le cas \struck{général} \add{où les $X_\alpha$ sont fermés}, un argument de descente (\uncertain{commode} à \uncertain{ignorer} dans la première hypothèse) donne le résultat énoncé.

Cette dernière démonstration — \uncertain{l'avantage} \add{\uncertain{entre autres}} de rendre également \uncertain{plausible} le résultat un peu plus général, et peut-être plus intéressant dans les pratiques où pour tout $\alpha$ on se donne \add{non vides}
\[
I_\alpha \subset T_\alpha \quad \text{partie de } T_\alpha \ (\text{\uncertain{non vide}, \uncertain{finie}}),
\]
partie finie) de telle façon que pour $T_\alpha \subset T_\beta$ on ait $I_\alpha \subset I_\beta$ (\uncertain{ou} plutôt, on peut définir $I \subset X_K$ convenable et \add{\uncertain{poser}} $I_\alpha = I \cap T_\alpha$, donc $I = \bigcup I_\alpha$), on se borne aux $g^{\alpha}_{y,x}$ pour $x, y \in I_\alpha$ : \uncertain{leur} connaissance permet de retrouver \emph{tous} \struck{\ill{}} $g_l(u)$, $l$ un chemin quelconque dans

Jusqu'à présent, on a maintenu $X$, $u$ fixés, on a fait varier $l$, resp. $(T/X, x, y)$. Si on regarde ce qui se passe pour $u \in \Gamma$ variable, la formule
\note{la phrase interrompue (« dans… ») reste en suspens ; un trait horizontal la sépare du paragraphe suivant, qui se poursuit sur la page suivante}

\marginal{note dense en oblique dans la marge gauche, sur la moitié inférieure de la page ; on y lit « $\Gamma$ », « $\pi_1(X,I)$ », « quelque soit », « $g_l(u)$ », « $f_l(u)$ », « $x, y \in I$ », « $T_\alpha \cap I$ », « chemin » ; le reste \ill{}}

\page{77}
essentielle triviale, qui fait pendant à (6), est
\[
\text{(28)}\qquad {}^{uv}l = {}^{u}({}^{v}l) \qquad u, v \in \Gamma,\ l : x \to y \ \text{chemin dans } \Pi_1(X)
\]
\note{le numéro (28) est surchargé ; il redouble celui de la page précédente}
\struck{d'où}, si $l$ est chemin de $\Pi_1(X)$, d'où
\[
{}^{v}l = g_l(v)\cdot l
\]
\struck{grâce à (6)} d'où on tire par (8)
\[
{}^{u}({}^{v}l) = {}^{u}\{g_l(v)\}\cdot\underbrace{{}^{u}l}_{g_l(u)\cdot l} = {}^{u}g_l(v)\,g_l(u)\cdot l
\]
d'où on tire
\[
\text{(29)}\qquad
\left|\begin{array}{l}
g_l(uv) = {}^{u}g_l(v)\cdot g_l(u) \quad\text{soit}\\
g_l^{-1}(uv) = g_l^{-1}(u)\cdot{}^{u}g_l^{-1}(v)
\end{array}\right.
\]
et de même
\[
\text{(30)}\qquad f_l(uv) = f_l(u)\cdot{}^{u}f_l(v)
\]
i.e. $f_l$ et $g_l^{-1}$ sont, pour $l$ fixé, $u$ variable, des 1-cocycles \add{sur $\Gamma$} à valeurs dans $\widehat{\pi}_1(X,x)$ resp. $\widehat{\pi}_1(X,y)$ (groupes sur lesquels $\Gamma$ opère), d'ailleurs déduits l'un de l'autre par l'isomorphisme
\[
\widehat{\pi}_1(X,x) \xrightarrow{\ \pi_1(l)\ } \widehat{\pi}_1(X,y) \quad \text{induit par } l \ldots
\]
Ici, c'est la formule pour $f_l$ qui a une gueule plus sympathique que celle pour $g_l$ (pour (10) (11) c'était l'inverse — \struck{\ill{}} \add{« \uncertain{multiplicativité} en $l$ »}). Les 1-cocycles obtenus (\struck{à \ill{}} \add{$f$} des 1-cocycles \add{droits}, $g$ des 1-cocycles à gauche) correspondent à des $\Gamma$-torseurs à droite resp. à gauche sous $\widehat{\pi}_1(X,x)$ resp.\ $\widehat{\pi}_1(X,y)$, qui bien sûr n'est autre que $\operatorname{Hom}_{\widehat{\Pi}_1(X)}(x,y)$
\note{la phrase se poursuit sur la page suivante}

\marginal{On pourrait expliquer \ill{} les \uncertain{deux} \ill{} $(f_l, g_l)$ \ill{} « bicocycle » \ill{} $\Gamma$ à valeurs dans $\bigl(\widehat{\pi}_1(X,y), \widehat{\pi}_1(X,x)\bigr)$ — note oblique dans la marge gauche, en face de (30), lecture incertaine}

\page{78}
\note{p. 7 de l'auteur}
dont $l$ est justement un élément jouant le rôle d'un « élément origine », pour en déduire des 1-cocycles…

\note{le paragraphe qui suit, jusqu'à « pour $x_0$ fixé », est barré de quatre traits obliques}
\struck{Le théorème de Mordell-Weil implique que si $K$ est un corps de type fini sur $\mathbb{Q}$, et si $X$ peut se plonger dans un \ill{}} \add{\struck{VA, \ill{}}} \struck{groupes algébriques ext. d'une VA par un tore, alors pour $x_0$ fixé,}

Si on garde $x, y$, et qu'on remplace $l$ par un $l' : x \to y$, alors \struck{les} on aura
\[
l' = la = bl \qquad a \in \pi_1(X,x),\ b \in \pi_1(X,y)
\]
et on aura, en vertu de (10) (11)
\[
\text{(31)}\qquad
\left\{
\begin{array}{l}
f_{l'}(u) = \operatorname{int}(a^{-1})\,f_l(u)\,\underbrace{f_a(u)}_{a^{-1}\cdot{}^{u}a} = a^{-1}\cdot f_l(u)\cdot{}^{u}a \\[1ex]
g_{l'}(u) = \underbrace{g_b(u)}_{{}^{u}b\cdot b^{-1}}\,\operatorname{int}(b)\,g_l(u) = {}^{u}b\,g_l(u)\,b^{-1}
\end{array}
\right.
\]
\[
\bigl(\text{i.e.}\ \ g^{-1}_{l'}(u) = b\,g^{-1}_l(u)\,{}^{u}b^{-1}\bigr)
\]
\note{sous le $g_l(u)$ de la deuxième ligne, une accolade avec une expression biffée}
i.e. $f_{l'}$ est un 1-cocycle \add{à droite} homologue à $f_l$, $g_{l'}$ un 1-cocycle homologue à $g_l$. Ainsi, pour $x, y \in X_K$, on trouve des éléments
\[
\text{(32)}\qquad
\left\{
\begin{array}{l}
g_{y,x} \in H^1\bigl(\Gamma, \widehat{\pi}_1(X,y)\bigr) \\
f_{y,x} \in H^1_{\mathrm{dr}}\bigl(\Gamma, \widehat{\pi}_1(X,x)\bigr)
\end{array}
\right.
\]
\note{le numéro (32) se lit dans la note marginale, non devant la formule ; sous le $H^1$ de la première ligne, un signe effacé}
\struck{\ill{}} qui sont « associés » en un sens convenable, via les \uncertain{bijections} canoniques
\[
\widehat{\pi}_1(X,x) \xrightarrow[\text{ext}]{\ \text{isom}\ } \widehat{\pi}_1(Y,y)
\]
\note{on lit $\widehat{\pi}_1(Y,y)$ ; on attendrait $\widehat{\pi}_1(X,y)$}

\marginal{Les $f_l(u)$, $g_l(u)$ se déduisent \ill{} \ill{} \ill{} $l'$ \ill{} $\widehat{\Pi}_1$ \ill{} $x, y \in X_K$ \ill{} de la formule (31) \ill{} $\widehat{\pi}_1(X)$. \ill{} définition des classes de \ill{} (32) \ill{} de \ill{} — longue note oblique dans la marge gauche, sur la moitié inférieure de la page ; presque tout \ill{}}

\page{79}
toute la structure étant décrite commodément en termes du bitorseur $\operatorname{Hom}_{\widehat{\Pi}_1(X)}(x,y)$ sous $\bigl(\widehat{\pi}_1(X,y), \widehat{\pi}_1(X,x)\bigr)$ et l'opération de $\Gamma$ sur tout ce fourbis. On se rappellera seulement que $f_{y,x}$ et $g_{y,x}$ se déterminent mutuellement, et qu'ils sont déterminés par la connaissance d'\emph{un seul} \struck{$f_l$ : $\Gamma \to \widehat{\pi}_1(X,x)$} \struck{(i.e. $f_l(u)$} $f_l$, ou $g_l$
\[
\text{(33)}\qquad
\left\{
\begin{array}{l}
f_l : \Gamma \longrightarrow \widehat{\pi}_1(X,x) \\
g_l : \Gamma \longrightarrow \widehat{\pi}_1(Y,y)
\end{array}
\right.
\qquad l : x \to y \ \ (\text{dans } \Pi_1(X))
\]
(si on a \struck{un} un $T^{\alpha}$ 1-connexe $\subset X$, contenant $x, y$) ils sont connus quand on connaît seulement $g^{T^{\alpha}}_{x,y}$.

\marginal{NB $f_{x,y}$ et $f_{y,x}$, \uncertain{et} $g_{x,y}$ et $g_{y,x}$, se déterminent mutuellement, p.ex. \ill{} $l^{-1}$ \ill{} (34) $g_{l^{-1}}(u) = f_l(u)^{-1}$, $f_{l^{-1}}(u) = g_l(u)^{-1}$ — note oblique dans la marge gauche, lecture incertaine}

Quand $X$ se plonge dans une VA, \uncertain{ou} plus généralement s'il se plonge dans une extension d'une VA par un tore (p.ex. $X$ courbe alg. « anabélienne »), alors le th. de Mordell-Weil implique que pour un $x_0 \in X_K$ fixé (\ill{} — mais en $x_0$ \uncertain{varié} pour une \ill{} bien !)
\note{la fin de la page est mal lisible ; la phrase principale se poursuit sur la page suivante (« un point $x \in X_K$ est entièrement déterminé… »)}

\page{80}
\note{p. 8 de l'auteur}
un point $x \in X_K$ est entièrement déterminé par la connaissance de $f_{x_0,x}$
\[
f_{x_0,x} \in H^1_{\mathrm{dr}}\bigl(\Gamma, \underbrace{\widehat{\pi}_1(X,x_0)}_{\widehat{\pi}(x_0)}\bigr),
\]
et a fortiori par la connaissance du \add{1-cocycle}
\[
f_l : \Gamma \longrightarrow \widehat{\pi}(x_0)
\]
pour un seul $l : x_0 \to x$ — ou ce qui revient au même, par la connaissance de $g^{-1}_{l^{-1}} \in H^1_g\bigl(\Gamma, \widehat{\pi}_1(x_0)\bigr)$, où $l^{-1} : x \to x_0$.
\marginal{$f_l = (g_{l^{-1}})^{-1}$}
\note{dans $f_l : \Gamma \to \widehat{\pi}(x_0)$, un indice est surchargé sur le $\pi$ ; dans $g^{-1}_{l^{-1}}$, l'indice est écrit sur un autre}

La détermination « explicite » d'une $f_l$ — \struck{d'une $g_l^{-1}$} (qui caractérise donc $x$ dans le cas envisagé) est un pb très intéressant, et sûrement profond. Comme \struck{\ill{}} le groupe $\Gamma$ lui-même est loin d'être connu et à la \uncertain{main}, il faut préciser d'ailleurs ce qu'on peut bien prendre pour une « détermination » (\uncertain{tout à fait} explicite !) de $f_l$. Pour ceci, il faut sûrement considérer que $u \in \Gamma$ est repéré par les invariants $f_{l_j}(u) \in \widehat{\pi}_1(x_j)$ ou $g_{l_j}(u) \in \widehat{\pi}_1(y_j)$, associés \add{aux \uncertain{lacets} $l_j : x_j \to y_j$} \add{\uncertain{d'un}} système \add{fini} de générateurs du groupoïde $\Pi_1(X, I)$ ($j \in J$, un indice pour les flèches génératrices, $x_j, y_j \in I$),
\note{la phrase, et l'argument, se poursuivent au-delà de ce lot}
\end{document}
